% --- Archon physics formalization source begin ---
% archon:physics
% archon:covers PhyXMiniProblems/problem_phyx_mini_0961.lean
% archon:source-report reports/phyx_mini/problem_phyx_mini_0961.source.json

\chapter{Physics problem phyx\_mini\_0961}
\label{ch:PhyXMiniProblems_problem_phyx_mini_0961}

\paragraph{Problem source.}
\#\# Physical scenario

A wire carrying a \textbackslash{}( 28.0\textbackslash{},\textbackslash{}text\{A\} \textbackslash{}) current bends through a right angle. Consider two \textbackslash{}( 2.00\textbackslash{},\textbackslash{}text\{mm\} \textbackslash{}) segments of wire, each \textbackslash{}( 3.00\textbackslash{},\textbackslash{}text\{cm\} \textbackslash{}) from the bend.

\#\# Question

Find the magnitude of the magnetic field these two segments produce at point \textbackslash{}( P \textbackslash{}), which is midway between them.

\#\# Answer choices

- A: A: 2.12 \textbackslash{}times 10\textasciicircum{}\{-5\}\textbackslash{}

- B: B: 1.76 \textbackslash{}times 10\textasciicircum{}\{-5\}\textbackslash{}

- C: C: 8.80 \textbackslash{}times 10\textasciicircum{}\{-6\}\textbackslash{}

- D: D: 1.76 \textbackslash{}times 10\textasciicircum{}\{-4\}\textbackslash{}

\#\# Figure caption (auxiliary; use the image as primary evidence)

The image depicts a diagram of a bent wire forming a right angle. 

- The wire has two straight segments, each labeled with their length as 3.00 cm. 
- The thickness of the wire is consistent, with a diameter of 2.00 mm on both segments.
- There is a purple arrow labeled "I" pointing upwards along the vertical segment, indicating the direction of current flow.
- A point labeled "P" is marked with a black dot, positioned diagonally from the bend, and shown with a dashed line connecting the dot to the bend. The length of this diagonal is not explicitly labeled.
- The vertical and horizontal segments are both marked with brackets showing their lengths as "3.00 cm".

\#\# Dataset metadata

Magnetism; Spatial Relation Reasoning, Physical Model Grounding Reasoning

\paragraph{Recorded answer/context.}
B: B: 1.76 \textbackslash{}times 10\textasciicircum{}\{-5\}\textbackslash{}

\paragraph{Figure/image path.}
/root/proposal\_for\_physic/science-mango/phyx\_mini\_run/phyx\_data/test\_image/961.png

\paragraph{Formalization target.}
create a compiling Lean file with sorry bodies at `PhyXMiniProblems/problem\_phyx\_mini\_0961.lean`. The Lean declarations must preserve the physical quantities, dimensions or dimensional roles, figure labels, governing-law hypotheses, and final relation expressed by this problem.
Use Mathlib/Physlib names found through LeanExplore where available. If a physics API is missing, introduce faithful local abstractions rather than scalar placeholder aliases.

\begin{theorem}[Physics formalization target]
\label{thm:physics:phyx_mini_0961:target}
\lean{PhyXMiniProblems.ProblemPhyXMini0961.magneticFieldMagnitudeAtMidpoint}
\uses{def:physics:phyx-mini-0961:phyxminiproblems-problemphyxmini0961-rightanglewiresetup, def:physics:phyx-mini-0961:phyxminiproblems-problemphyxmini0961-totalfieldmagnitudeatpinteslas, def:physics:phyx-mini-0961:phyxminiproblems-problemphyxmini0961-matchesrightanglewireproblemdata, def:physics:phyx-mini-0961:phyxminiproblems-problemphyxmini0961-matchesprimaryrightanglewirefigure, def:physics:phyx-mini-0961:phyxminiproblems-problemphyxmini0961-hasphysicalrightanglewireparameters, def:physics:phyx-mini-0961:phyxminiproblems-problemphyxmini0961-usesstandardvacuumpermeability, def:physics:phyx-mini-0961:phyxminiproblems-problemphyxmini0961-satisfiessteadycurrentfieldmodel, def:physics:phyx-mini-0961:phyxminiproblems-problemphyxmini0961-satisfiesshortelementbiotsavartlaw, def:physics:phyx-mini-0961:phyxminiproblems-problemphyxmini0961-satisfiestwoelementsuperposition, def:physics:phyx-mini-0961:phyxminiproblems-problemphyxmini0961-magnitudereadoutmatchestotalfield, def:physics:phyx-mini-0961:phyxminiproblems-problemphyxmini0961-displayedmagneticfieldinteslas, def:physics:phyx-mini-0961:phyxminiproblems-problemphyxmini0961-roundstochoiceprecision, def:physics:phyx-mini-0961:phyxminiproblems-problemphyxmini0961-isuniqueclosestdisplayedanswer}
The assigned autoformalize agent should translate this physics problem into Lean declarations in the covered file, with theorem and lemma proof bodies written as `by sorry`.
\end{theorem}
\begin{proof}
This is an autoformalization task, not a proof task. Produce statements that can later be proved by the physics prover without weakening the physical model.
\end{proof}

% --- Archon declaration topology begin ---
\section{Declaration topology}
\begin{definition}[electric Current Dimension]
  \label{def:physics:phyx-mini-0961:phyxminiproblems-problemphyxmini0961-electriccurrentdimension}
  \lean{PhyXMiniProblems.ProblemPhyXMini0961.electricCurrentDimension}
  Electric current has dimension charge per unit time
\end{definition}
\begin{proof}
  This is the declared model definition of electric Current Dimension.
\end{proof}
\begin{definition}[magnetic Permeability Dimension]
  \label{def:physics:phyx-mini-0961:phyxminiproblems-problemphyxmini0961-magneticpermeabilitydimension}
  \lean{PhyXMiniProblems.ProblemPhyXMini0961.magneticPermeabilityDimension}
  Magnetic permeability has dimension M L C⁻²
\end{definition}
\begin{proof}
  This is the declared model definition of magnetic Permeability Dimension.
\end{proof}
\begin{definition}[magnetic Flux Density Dimension]
  \label{def:physics:phyx-mini-0961:phyxminiproblems-problemphyxmini0961-magneticfluxdensitydimension}
  \lean{PhyXMiniProblems.ProblemPhyXMini0961.magneticFluxDensityDimension}
  Magnetic flux density (tesla) has dimension M T⁻¹ C⁻¹
\end{definition}
\begin{proof}
  This is the declared model definition of magnetic Flux Density Dimension.
\end{proof}
\begin{definition}[Length Quantity]
  \label{def:physics:phyx-mini-0961:phyxminiproblems-problemphyxmini0961-lengthquantity}
  \lean{PhyXMiniProblems.ProblemPhyXMini0961.LengthQuantity}
  A nonnegative, unit-independent physical length
\end{definition}
\begin{proof}
  This is the declared model definition of Length Quantity.
\end{proof}
\begin{definition}[Electric Current Magnitude]
  \label{def:physics:phyx-mini-0961:phyxminiproblems-problemphyxmini0961-electriccurrentmagnitude}
  \lean{PhyXMiniProblems.ProblemPhyXMini0961.ElectricCurrentMagnitude}
  \uses{def:physics:phyx-mini-0961:phyxminiproblems-problemphyxmini0961-electriccurrentdimension}
  A nonnegative, unit-independent electric-current magnitude
\end{definition}
\begin{proof}
  This is the declared model definition of Electric Current Magnitude.
\end{proof}
\begin{definition}[Magnetic Permeability Quantity]
  \label{def:physics:phyx-mini-0961:phyxminiproblems-problemphyxmini0961-magneticpermeabilityquantity}
  \lean{PhyXMiniProblems.ProblemPhyXMini0961.MagneticPermeabilityQuantity}
  \uses{def:physics:phyx-mini-0961:phyxminiproblems-problemphyxmini0961-magneticpermeabilitydimension}
  A nonnegative, unit-independent magnetic permeability
\end{definition}
\begin{proof}
  This is the declared model definition of Magnetic Permeability Quantity.
\end{proof}
\begin{definition}[Magnetic Flux Density Magnitude]
  \label{def:physics:phyx-mini-0961:phyxminiproblems-problemphyxmini0961-magneticfluxdensitymagnitude}
  \lean{PhyXMiniProblems.ProblemPhyXMini0961.MagneticFluxDensityMagnitude}
  \uses{def:physics:phyx-mini-0961:phyxminiproblems-problemphyxmini0961-magneticfluxdensitydimension}
  A nonnegative, unit-independent magnetic-flux-density magnitude
\end{definition}
\begin{proof}
  This is the declared model definition of Magnetic Flux Density Magnitude.
\end{proof}
\begin{definition}[coherent SIReadout]
  \label{def:physics:phyx-mini-0961:phyxminiproblems-problemphyxmini0961-coherentsireadout}
  \lean{PhyXMiniProblems.ProblemPhyXMini0961.coherentSIReadout}
  Read a nonnegative physical quantity in coherent SI units
\end{definition}
\begin{proof}
  This is the declared model definition of coherent SIReadout.
\end{proof}
\begin{definition}[length In Meters]
  \label{def:physics:phyx-mini-0961:phyxminiproblems-problemphyxmini0961-lengthinmeters}
  \lean{PhyXMiniProblems.ProblemPhyXMini0961.lengthInMeters}
  \uses{def:physics:phyx-mini-0961:phyxminiproblems-problemphyxmini0961-lengthquantity, def:physics:phyx-mini-0961:phyxminiproblems-problemphyxmini0961-coherentsireadout}
  Read a physical length in metres
\end{definition}
\begin{proof}
  This is the declared model definition of length In Meters.
\end{proof}
\begin{definition}[length In Centimeters]
  \label{def:physics:phyx-mini-0961:phyxminiproblems-problemphyxmini0961-lengthincentimeters}
  \lean{PhyXMiniProblems.ProblemPhyXMini0961.lengthInCentimeters}
  \uses{def:physics:phyx-mini-0961:phyxminiproblems-problemphyxmini0961-lengthquantity, def:physics:phyx-mini-0961:phyxminiproblems-problemphyxmini0961-lengthinmeters}
  Read a physical length in centimetres
\end{definition}
\begin{proof}
  This is the declared model definition of length In Centimeters.
\end{proof}
\begin{definition}[length In Millimeters]
  \label{def:physics:phyx-mini-0961:phyxminiproblems-problemphyxmini0961-lengthinmillimeters}
  \lean{PhyXMiniProblems.ProblemPhyXMini0961.lengthInMillimeters}
  \uses{def:physics:phyx-mini-0961:phyxminiproblems-problemphyxmini0961-lengthquantity, def:physics:phyx-mini-0961:phyxminiproblems-problemphyxmini0961-lengthinmeters}
  Read a physical length in millimetres
\end{definition}
\begin{proof}
  This is the declared model definition of length In Millimeters.
\end{proof}
\begin{definition}[current In Amperes]
  \label{def:physics:phyx-mini-0961:phyxminiproblems-problemphyxmini0961-currentinamperes}
  \lean{PhyXMiniProblems.ProblemPhyXMini0961.currentInAmperes}
  \uses{def:physics:phyx-mini-0961:phyxminiproblems-problemphyxmini0961-electriccurrentmagnitude, def:physics:phyx-mini-0961:phyxminiproblems-problemphyxmini0961-coherentsireadout}
  Read a physical current magnitude in amperes
\end{definition}
\begin{proof}
  This is the declared model definition of current In Amperes.
\end{proof}
\begin{definition}[permeability In Tesla Meters Per Ampere]
  \label{def:physics:phyx-mini-0961:phyxminiproblems-problemphyxmini0961-permeabilityinteslametersperampere}
  \lean{PhyXMiniProblems.ProblemPhyXMini0961.permeabilityInTeslaMetersPerAmpere}
  \uses{def:physics:phyx-mini-0961:phyxminiproblems-problemphyxmini0961-magneticpermeabilityquantity, def:physics:phyx-mini-0961:phyxminiproblems-problemphyxmini0961-coherentsireadout}
  Read magnetic permeability in tesla-metres per ampere
\end{definition}
\begin{proof}
  This is the declared model definition of permeability In Tesla Meters Per Ampere.
\end{proof}
\begin{definition}[magnetic Flux Density In Teslas]
  \label{def:physics:phyx-mini-0961:phyxminiproblems-problemphyxmini0961-magneticfluxdensityinteslas}
  \lean{PhyXMiniProblems.ProblemPhyXMini0961.magneticFluxDensityInTeslas}
  \uses{def:physics:phyx-mini-0961:phyxminiproblems-problemphyxmini0961-magneticfluxdensitymagnitude, def:physics:phyx-mini-0961:phyxminiproblems-problemphyxmini0961-coherentsireadout}
  Read a magnetic-flux-density magnitude in teslas
\end{definition}
\begin{proof}
  This is the declared model definition of magnetic Flux Density In Teslas.
\end{proof}
\begin{definition}[Spatial Vector]
  \label{def:physics:phyx-mini-0961:phyxminiproblems-problemphyxmini0961-spatialvector}
  \lean{PhyXMiniProblems.ProblemPhyXMini0961.SpatialVector}
  Three-dimensional coherent-SI spatial vectors
\end{definition}
\begin{proof}
  This is the declared model definition of Spatial Vector.
\end{proof}
\begin{definition}[spatial Cross Product]
  \label{def:physics:phyx-mini-0961:phyxminiproblems-problemphyxmini0961-spatialcrossproduct}
  \lean{PhyXMiniProblems.ProblemPhyXMini0961.spatialCrossProduct}
  \uses{def:physics:phyx-mini-0961:phyxminiproblems-problemphyxmini0961-spatialvector}
  Mathlib's right-handed cross product transported to SpatialVector
\end{definition}
\begin{proof}
  This is the declared model definition of spatial Cross Product.
\end{proof}
\begin{definition}[Coordinate Axis]
  \label{def:physics:phyx-mini-0961:phyxminiproblems-problemphyxmini0961-coordinateaxis}
  \lean{PhyXMiniProblems.ProblemPhyXMini0961.CoordinateAxis}
  Coordinate axes used to orient the two wire legs and the page normal
\end{definition}
\begin{proof}
  This is the declared model definition of Coordinate Axis.
\end{proof}
\begin{definition}[axis Vector]
  \label{def:physics:phyx-mini-0961:phyxminiproblems-problemphyxmini0961-axisvector}
  \lean{PhyXMiniProblems.ProblemPhyXMini0961.axisVector}
  \uses{def:physics:phyx-mini-0961:phyxminiproblems-problemphyxmini0961-spatialvector, def:physics:phyx-mini-0961:phyxminiproblems-problemphyxmini0961-coordinateaxis}
  Positive unit vector associated with a Cartesian axis
\end{definition}
\begin{proof}
  This is the declared model definition of axis Vector.
\end{proof}
\begin{definition}[Wire Leg]
  \label{def:physics:phyx-mini-0961:phyxminiproblems-problemphyxmini0961-wireleg}
  \lean{PhyXMiniProblems.ProblemPhyXMini0961.WireLeg}
  The two perpendicular legs of the bent wire
\end{definition}
\begin{proof}
  This is the declared model definition of Wire Leg.
\end{proof}
\begin{definition}[Marked Element]
  \label{def:physics:phyx-mini-0961:phyxminiproblems-problemphyxmini0961-markedelement}
  \lean{PhyXMiniProblems.ProblemPhyXMini0961.MarkedElement}
  The marked short current element on each leg
\end{definition}
\begin{proof}
  This is the declared model definition of Marked Element.
\end{proof}
\begin{definition}[leg]
  \label{def:physics:phyx-mini-0961:phyxminiproblems-problemphyxmini0961-markedelement-leg}
  \lean{PhyXMiniProblems.ProblemPhyXMini0961.MarkedElement.leg}
  \uses{def:physics:phyx-mini-0961:phyxminiproblems-problemphyxmini0961-wireleg, def:physics:phyx-mini-0961:phyxminiproblems-problemphyxmini0961-markedelement}
  The leg containing a marked current element
\end{definition}
\begin{proof}
  This is the declared model definition of leg.
\end{proof}
\begin{definition}[Axis Direction]
  \label{def:physics:phyx-mini-0961:phyxminiproblems-problemphyxmini0961-axisdirection}
  \lean{PhyXMiniProblems.ProblemPhyXMini0961.AxisDirection}
  Qualitative directions used by the current arrow and wire legs
\end{definition}
\begin{proof}
  This is the declared model definition of Axis Direction.
\end{proof}
\begin{definition}[direction Vector]
  \label{def:physics:phyx-mini-0961:phyxminiproblems-problemphyxmini0961-directionvector}
  \lean{PhyXMiniProblems.ProblemPhyXMini0961.directionVector}
  \uses{def:physics:phyx-mini-0961:phyxminiproblems-problemphyxmini0961-spatialvector, def:physics:phyx-mini-0961:phyxminiproblems-problemphyxmini0961-axisvector, def:physics:phyx-mini-0961:phyxminiproblems-problemphyxmini0961-axisdirection}
  Unit vector represented by a qualitative direction
\end{definition}
\begin{proof}
  This is the declared model definition of direction Vector.
\end{proof}
\begin{definition}[Current Regime]
  \label{def:physics:phyx-mini-0961:phyxminiproblems-problemphyxmini0961-currentregime}
  \lean{PhyXMiniProblems.ProblemPhyXMini0961.CurrentRegime}
  A time-independent or time-varying current regime
\end{definition}
\begin{proof}
  This is the declared model definition of Current Regime.
\end{proof}
\begin{definition}[Right Angle Wire Figure]
  \label{def:physics:phyx-mini-0961:phyxminiproblems-problemphyxmini0961-rightanglewirefigure}
  \lean{PhyXMiniProblems.ProblemPhyXMini0961.RightAngleWireFigure}
  \uses{def:physics:phyx-mini-0961:phyxminiproblems-problemphyxmini0961-wireleg, def:physics:phyx-mini-0961:phyxminiproblems-problemphyxmini0961-markedelement, def:physics:phyx-mini-0961:phyxminiproblems-problemphyxmini0961-axisdirection}
  Literal graphical readouts from the primary image 961.png. The scalar labels remain distinct from the dimensionful quantities they calibrate
\end{definition}
\begin{proof}
  This is the declared model definition of Right Angle Wire Figure.
\end{proof}
\begin{definition}[Right Angle Wire Setup]
  \label{def:physics:phyx-mini-0961:phyxminiproblems-problemphyxmini0961-rightanglewiresetup}
  \lean{PhyXMiniProblems.ProblemPhyXMini0961.RightAngleWireSetup}
  \uses{def:physics:phyx-mini-0961:phyxminiproblems-problemphyxmini0961-lengthquantity, def:physics:phyx-mini-0961:phyxminiproblems-problemphyxmini0961-electriccurrentmagnitude, def:physics:phyx-mini-0961:phyxminiproblems-problemphyxmini0961-magneticpermeabilityquantity, def:physics:phyx-mini-0961:phyxminiproblems-problemphyxmini0961-magneticfluxdensitymagnitude, def:physics:phyx-mini-0961:phyxminiproblems-problemphyxmini0961-spatialvector, def:physics:phyx-mini-0961:phyxminiproblems-problemphyxmini0961-markedelement, def:physics:phyx-mini-0961:phyxminiproblems-problemphyxmini0961-currentregime, def:physics:phyx-mini-0961:phyxminiproblems-problemphyxmini0961-rightanglewirefigure}
  Independent physical quantities and observables for the bent wire. The requested magnetic-field magnitude is an observable, not a definition from an answer choice or target numeral
\end{definition}
\begin{proof}
  This is the declared model definition of Right Angle Wire Setup.
\end{proof}
\begin{definition}[displacement To PIn Meters]
  \label{def:physics:phyx-mini-0961:phyxminiproblems-problemphyxmini0961-displacementtopinmeters}
  \lean{PhyXMiniProblems.ProblemPhyXMini0961.displacementToPInMeters}
  \uses{def:physics:phyx-mini-0961:phyxminiproblems-problemphyxmini0961-spatialvector, def:physics:phyx-mini-0961:phyxminiproblems-problemphyxmini0961-markedelement, def:physics:phyx-mini-0961:phyxminiproblems-problemphyxmini0961-rightanglewiresetup}
  Displacement from a marked element's center to P, in metres
\end{definition}
\begin{proof}
  This is the declared model definition of displacement To PIn Meters.
\end{proof}
\begin{definition}[element Field Vector At PIn Teslas]
  \label{def:physics:phyx-mini-0961:phyxminiproblems-problemphyxmini0961-elementfieldvectoratpinteslas}
  \lean{PhyXMiniProblems.ProblemPhyXMini0961.elementFieldVectorAtPInTeslas}
  \uses{def:physics:phyx-mini-0961:phyxminiproblems-problemphyxmini0961-spatialvector, def:physics:phyx-mini-0961:phyxminiproblems-problemphyxmini0961-markedelement, def:physics:phyx-mini-0961:phyxminiproblems-problemphyxmini0961-rightanglewiresetup}
  Magnetic-field vector at P due to one marked element, in teslas
\end{definition}
\begin{proof}
  This is the declared model definition of element Field Vector At PIn Teslas.
\end{proof}
\begin{definition}[total Field Vector At PIn Teslas]
  \label{def:physics:phyx-mini-0961:phyxminiproblems-problemphyxmini0961-totalfieldvectoratpinteslas}
  \lean{PhyXMiniProblems.ProblemPhyXMini0961.totalFieldVectorAtPInTeslas}
  \uses{def:physics:phyx-mini-0961:phyxminiproblems-problemphyxmini0961-spatialvector, def:physics:phyx-mini-0961:phyxminiproblems-problemphyxmini0961-rightanglewiresetup}
  Total magnetic-field vector at P, in teslas
\end{definition}
\begin{proof}
  This is the declared model definition of total Field Vector At PIn Teslas.
\end{proof}
\begin{definition}[total Field Magnitude At PIn Teslas]
  \label{def:physics:phyx-mini-0961:phyxminiproblems-problemphyxmini0961-totalfieldmagnitudeatpinteslas}
  \lean{PhyXMiniProblems.ProblemPhyXMini0961.totalFieldMagnitudeAtPInTeslas}
  \uses{def:physics:phyx-mini-0961:phyxminiproblems-problemphyxmini0961-magneticfluxdensityinteslas, def:physics:phyx-mini-0961:phyxminiproblems-problemphyxmini0961-rightanglewiresetup}
  Physical magnitude of the total field at P, read in teslas
\end{definition}
\begin{proof}
  This is the declared model definition of total Field Magnitude At PIn Teslas.
\end{proof}
\begin{definition}[Matches Right Angle Wire Problem Data]
  \label{def:physics:phyx-mini-0961:phyxminiproblems-problemphyxmini0961-matchesrightanglewireproblemdata}
  \lean{PhyXMiniProblems.ProblemPhyXMini0961.MatchesRightAngleWireProblemData}
  \uses{def:physics:phyx-mini-0961:phyxminiproblems-problemphyxmini0961-lengthinmeters, def:physics:phyx-mini-0961:phyxminiproblems-problemphyxmini0961-lengthincentimeters, def:physics:phyx-mini-0961:phyxminiproblems-problemphyxmini0961-lengthinmillimeters, def:physics:phyx-mini-0961:phyxminiproblems-problemphyxmini0961-currentinamperes, def:physics:phyx-mini-0961:phyxminiproblems-problemphyxmini0961-rightanglewiresetup}
  Numerical problem data, separated from the requested field value
\end{definition}
\begin{proof}
  This is the declared model definition of Matches Right Angle Wire Problem Data.
\end{proof}
\begin{definition}[Matches Primary Right Angle Wire Figure]
  \label{def:physics:phyx-mini-0961:phyxminiproblems-problemphyxmini0961-matchesprimaryrightanglewirefigure}
  \lean{PhyXMiniProblems.ProblemPhyXMini0961.MatchesPrimaryRightAngleWireFigure}
  \uses{def:physics:phyx-mini-0961:phyxminiproblems-problemphyxmini0961-lengthinmeters, def:physics:phyx-mini-0961:phyxminiproblems-problemphyxmini0961-lengthincentimeters, def:physics:phyx-mini-0961:phyxminiproblems-problemphyxmini0961-lengthinmillimeters, def:physics:phyx-mini-0961:phyxminiproblems-problemphyxmini0961-currentinamperes, def:physics:phyx-mini-0961:phyxminiproblems-problemphyxmini0961-axisvector, def:physics:phyx-mini-0961:phyxminiproblems-problemphyxmini0961-directionvector, def:physics:phyx-mini-0961:phyxminiproblems-problemphyxmini0961-rightanglewiresetup}
  Primary-image transcription and coherent-SI coordinate geometry. The vertical element lies below the bend and carries current upward; after the bend, the horizontal element carries the same current rightward. P is the affine midpoint of the two element centers
\end{definition}
\begin{proof}
  This is the declared model definition of Matches Primary Right Angle Wire Figure.
\end{proof}
\begin{definition}[Has Physical Right Angle Wire Parameters]
  \label{def:physics:phyx-mini-0961:phyxminiproblems-problemphyxmini0961-hasphysicalrightanglewireparameters}
  \lean{PhyXMiniProblems.ProblemPhyXMini0961.HasPhysicalRightAngleWireParameters}
  \uses{def:physics:phyx-mini-0961:phyxminiproblems-problemphyxmini0961-lengthinmeters, def:physics:phyx-mini-0961:phyxminiproblems-problemphyxmini0961-currentinamperes, def:physics:phyx-mini-0961:phyxminiproblems-problemphyxmini0961-permeabilityinteslametersperampere, def:physics:phyx-mini-0961:phyxminiproblems-problemphyxmini0961-rightanglewiresetup, def:physics:phyx-mini-0961:phyxminiproblems-problemphyxmini0961-displacementtopinmeters}
  Positivity conditions selecting the physical, nondegenerate branch
\end{definition}
\begin{proof}
  This is the declared model definition of Has Physical Right Angle Wire Parameters.
\end{proof}
\begin{definition}[Uses Standard Vacuum Permeability]
  \label{def:physics:phyx-mini-0961:phyxminiproblems-problemphyxmini0961-usesstandardvacuumpermeability}
  \lean{PhyXMiniProblems.ProblemPhyXMini0961.UsesStandardVacuumPermeability}
  \uses{def:physics:phyx-mini-0961:phyxminiproblems-problemphyxmini0961-permeabilityinteslametersperampere, def:physics:phyx-mini-0961:phyxminiproblems-problemphyxmini0961-rightanglewiresetup}
  The dimensionful permeability agrees with the Physlib electromagnetic system and has the standard coherent-SI value 4π * 10⁻⁷ T m/A
\end{definition}
\begin{proof}
  This is the declared model definition of Uses Standard Vacuum Permeability.
\end{proof}
\begin{definition}[Satisfies Steady Current Field Model]
  \label{def:physics:phyx-mini-0961:phyxminiproblems-problemphyxmini0961-satisfiessteadycurrentfieldmodel}
  \lean{PhyXMiniProblems.ProblemPhyXMini0961.SatisfiesSteadyCurrentFieldModel}
  \uses{def:physics:phyx-mini-0961:phyxminiproblems-problemphyxmini0961-rightanglewiresetup}
  A steady current produces time-independent element and total fields
\end{definition}
\begin{proof}
  This is the declared model definition of Satisfies Steady Current Field Model.
\end{proof}
\begin{definition}[Satisfies Short Element Biot Savart Law]
  \label{def:physics:phyx-mini-0961:phyxminiproblems-problemphyxmini0961-satisfiesshortelementbiotsavartlaw}
  \lean{PhyXMiniProblems.ProblemPhyXMini0961.SatisfiesShortElementBiotSavartLaw}
  \uses{def:physics:phyx-mini-0961:phyxminiproblems-problemphyxmini0961-currentinamperes, def:physics:phyx-mini-0961:phyxminiproblems-problemphyxmini0961-permeabilityinteslametersperampere, def:physics:phyx-mini-0961:phyxminiproblems-problemphyxmini0961-spatialcrossproduct, def:physics:phyx-mini-0961:phyxminiproblems-problemphyxmini0961-rightanglewiresetup, def:physics:phyx-mini-0961:phyxminiproblems-problemphyxmini0961-displacementtopinmeters, def:physics:phyx-mini-0961:phyxminiproblems-problemphyxmini0961-elementfieldvectoratpinteslas}
  Short-current-element Biot--Savart law at P for both marked elements: B = (μ₀ I / (4π |r|³)) (dℓ × r). This general vector relation retains the separation, angle factor, and right-hand-rule direction. It contains no requested numerical field value
\end{definition}
\begin{proof}
  This is the declared model definition of Satisfies Short Element Biot Savart Law.
\end{proof}
\begin{definition}[Satisfies Two Element Superposition]
  \label{def:physics:phyx-mini-0961:phyxminiproblems-problemphyxmini0961-satisfiestwoelementsuperposition}
  \lean{PhyXMiniProblems.ProblemPhyXMini0961.SatisfiesTwoElementSuperposition}
  \uses{def:physics:phyx-mini-0961:phyxminiproblems-problemphyxmini0961-rightanglewiresetup, def:physics:phyx-mini-0961:phyxminiproblems-problemphyxmini0961-elementfieldvectoratpinteslas, def:physics:phyx-mini-0961:phyxminiproblems-problemphyxmini0961-totalfieldvectoratpinteslas}
  Magnetic fields produced by the two current elements superpose linearly
\end{definition}
\begin{proof}
  This is the declared model definition of Satisfies Two Element Superposition.
\end{proof}
\begin{definition}[Magnitude Readout Matches Total Field]
  \label{def:physics:phyx-mini-0961:phyxminiproblems-problemphyxmini0961-magnitudereadoutmatchestotalfield}
  \lean{PhyXMiniProblems.ProblemPhyXMini0961.MagnitudeReadoutMatchesTotalField}
  \uses{def:physics:phyx-mini-0961:phyxminiproblems-problemphyxmini0961-rightanglewiresetup, def:physics:phyx-mini-0961:phyxminiproblems-problemphyxmini0961-totalfieldvectoratpinteslas, def:physics:phyx-mini-0961:phyxminiproblems-problemphyxmini0961-totalfieldmagnitudeatpinteslas}
  The scalar physical readout is the Euclidean norm of the total field
\end{definition}
\begin{proof}
  This is the declared model definition of Magnitude Readout Matches Total Field.
\end{proof}
\begin{lemma}[total Field Magnitude At P exact]
  \label{lem:physics:phyx-mini-0961:phyxminiproblems-problemphyxmini0961-totalfieldmagnitudeatp-exact}
  \lean{PhyXMiniProblems.ProblemPhyXMini0961.totalFieldMagnitudeAtP_exact}
  \uses{def:physics:phyx-mini-0961:phyxminiproblems-problemphyxmini0961-rightanglewiresetup, def:physics:phyx-mini-0961:phyxminiproblems-problemphyxmini0961-totalfieldmagnitudeatpinteslas, def:physics:phyx-mini-0961:phyxminiproblems-problemphyxmini0961-matchesrightanglewireproblemdata, def:physics:phyx-mini-0961:phyxminiproblems-problemphyxmini0961-matchesprimaryrightanglewirefigure, def:physics:phyx-mini-0961:phyxminiproblems-problemphyxmini0961-hasphysicalrightanglewireparameters, def:physics:phyx-mini-0961:phyxminiproblems-problemphyxmini0961-usesstandardvacuumpermeability, def:physics:phyx-mini-0961:phyxminiproblems-problemphyxmini0961-satisfiessteadycurrentfieldmodel, def:physics:phyx-mini-0961:phyxminiproblems-problemphyxmini0961-satisfiesshortelementbiotsavartlaw, def:physics:phyx-mini-0961:phyxminiproblems-problemphyxmini0961-satisfiestwoelementsuperposition, def:physics:phyx-mini-0961:phyxminiproblems-problemphyxmini0961-magnitudereadoutmatchestotalfield}
  The two element fields have the same page-normal direction. The exact short-element model gives 14 * sqrt 2 / 1125000 T; this is approximately 1.75964 * 10⁻⁵ T
\end{lemma}
\begin{proof}
  The conclusion follows from the governing relations and figure readouts named in the statement of total Field Magnitude At P exact.
\end{proof}
\begin{definition}[Answer Choice]
  \label{def:physics:phyx-mini-0961:phyxminiproblems-problemphyxmini0961-answerchoice}
  \lean{PhyXMiniProblems.ProblemPhyXMini0961.AnswerChoice}
  Labels of the four magnetic-field-magnitude choices
\end{definition}
\begin{proof}
  This is the declared model definition of Answer Choice.
\end{proof}
\begin{definition}[displayed Magnetic Field In Teslas]
  \label{def:physics:phyx-mini-0961:phyxminiproblems-problemphyxmini0961-displayedmagneticfieldinteslas}
  \lean{PhyXMiniProblems.ProblemPhyXMini0961.displayedMagneticFieldInTeslas}
  \uses{def:physics:phyx-mini-0961:phyxminiproblems-problemphyxmini0961-answerchoice}
  Magnetic-field magnitude in teslas displayed beside each answer choice
\end{definition}
\begin{proof}
  This is the declared model definition of displayed Magnetic Field In Teslas.
\end{proof}
\begin{definition}[Rounds To Choice Precision]
  \label{def:physics:phyx-mini-0961:phyxminiproblems-problemphyxmini0961-roundstochoiceprecision}
  \lean{PhyXMiniProblems.ProblemPhyXMini0961.RoundsToChoicePrecision}
  Rounding to the nearest 0.01 * 10⁻⁵ T, the precision of the recorded choice 1.76 * 10⁻⁵ T
\end{definition}
\begin{proof}
  This is the declared model definition of Rounds To Choice Precision.
\end{proof}
\begin{definition}[Is Unique Closest Displayed Answer]
  \label{def:physics:phyx-mini-0961:phyxminiproblems-problemphyxmini0961-isuniqueclosestdisplayedanswer}
  \lean{PhyXMiniProblems.ProblemPhyXMini0961.IsUniqueClosestDisplayedAnswer}
  \uses{def:physics:phyx-mini-0961:phyxminiproblems-problemphyxmini0961-answerchoice, def:physics:phyx-mini-0961:phyxminiproblems-problemphyxmini0961-displayedmagneticfieldinteslas}
  A displayed choice is strictly closer than every distinct alternative
\end{definition}
\begin{proof}
  This is the declared model definition of Is Unique Closest Displayed Answer.
\end{proof}
% --- Archon declaration topology end ---
% --- Archon physics formalization source end ---
